%% Wiley 期刊通用投稿骨架
%% 说明:
%% - Wiley 正式模板（如 wileyNJD-v2）请在投稿时从 Wiley Author Services 下载对应期刊 class
%% - 本骨架基于标准 article，按 Wiley 常见结构组织
%% 编译: pdflatex 或 xelatex
\documentclass[11pt,a4paper]{article}

\usepackage[margin=1in]{geometry}
\usepackage{setspace}
\usepackage{graphicx}
\usepackage{amsmath,amssymb}
\usepackage{booktabs}
\usepackage{xcolor}
\usepackage{hyperref}

\hypersetup{
  colorlinks=true,
  linkcolor=black,
  citecolor=blue!50!black,
  urlcolor=blue!50!black,
  pdfauthor={Author Name},
  pdftitle={Wiley Journal Submission Template}
}

\setstretch{1.15}
\DeclareMathOperator*{\argmin}{arg\\,min}

\begin{document}

\begin{center}
{\LARGE\bfseries A Concise Template for Wiley Journal Submission}\\[1.2em]
{\large First A. Author$^{1}$, Second B. Author$^{1,2}$}\\[0.6em]
{$^{1}$ Department of Computer Science, University Name, City, Country}\\
{$^{2}$ Research Institute, Organization Name, City, Country}\\[0.6em]
\	exttt{first.author@example.com}
\end{center}

\vspace{1.2em}

\
oindent\	extbf{Abstract.}
This skeleton follows a structure commonly used for Wiley journal manuscripts. Replace the sample text with your abstract (typically 150--250 words). Provide keywords separately below the abstract.

\vspace{0.8em}

\
oindent\	extbf{Keywords:} LaTeX; Wiley; journal submission; template; scientific writing

\vspace{1.2em}

\section{Introduction}
Introduce the research problem, motivation, and contribution. Cite related work as required by the target journal.

\section{Related Work}
Summarize prior approaches and position your contribution relative to them.

\section{Methods}
Describe the framework, algorithms, or experimental setup.

\begin{equation}
\	heta^{*} = \argmin_{\	heta} \\, \mathcal{L}(\	heta; \mathcal{D})
\label{eq:objective}
\end{equation}

Equation~\eqref{eq:objective} is a placeholder objective function.

\section{Results}
Present quantitative findings with self-contained tables and figures.

\begin{table}[htbp]
\centering
\caption{Sample experimental comparison}
\label{tab:wiley-sample}
\begin{tabular}{lcc}
\	oprule
Model & Accuracy & F1 \\
\midrule
Baseline & 0.88 & 0.86 \\
Ours & 0.92 & 0.91 \\
\bottomrule
\end{tabular}
\end{table}

\section{Discussion}
Interpret the results, acknowledge limitations, and outline future work.

\section{Conclusion}
Restate the problem, method, and main finding in a short closing paragraph.

\section*{Acknowledgements}
We thank anonymous reviewers and colleagues for constructive feedback.

\section*{Conflict of interest}
The authors declare no conflict of interest.

\section*{Data availability statement}
Supporting data can be requested from the corresponding author.

\begin{thebibliography}{99}
\bibitem{ref1}
Author A. Title of the article. \	extit{Journal Name}. Year;Volume(Issue):Pages.

\bibitem{ref2}
Author B, Author C. Title of the paper. In: \	extit{Proceedings of the Conference}; Year; Pages.

\bibitem{ref3}
Author D. \	extit{Book Title}. Publisher; Year.
\end{thebibliography}

\end{document}
